\section{Évaluation constitutive et réalisation dimensionnelle}
\label{iii:sec:evaluacion}

L’évaluation réunit maintenant les constructions précédentes dans leur ordre effectif.
Les développements régionaux de $\pi$, $\varphi$ et $e$ proviennent du raffinement
nonadique ; le registre $K$ conserve ses douze positions. Sa lecture périodique
détermine la section positionnelle de $\alpha$. Sur celle-ci agissent le caractère
d’action et le retour régional, puis la paire angulaire et, enfin, la
réponse constitutive. Cette évaluation est postérieure à la sélection des
règles : aucun de ses chiffres n’est utilisé pour choisir un canal.

\subsection{Évaluation de la chaîne interne}

Pour reproduire le tableau, on emploie les développements régionaux à mille
décimales engendrés par le certificat nonadique et l’entier $N_K$ de
\eqref{eq:k-numerador}. La lecture utilisée est
\[
 \alpha=\pi+e-\varphi-4-\frac{N_K}{1000^{12}-1}.
\]
La section finie de douze blocs et la section prolongée ont leurs
définitions et leur différence exacte dans le chapitre du registre. Ici, on
maintient la seconde fixe pendant toute l’évaluation. On n’échange pas
les lectures lorsque la précision augmente, et on ne remplace pas l’opérateur analytique
de la seconde voie par une coïncidence décimale avec la première.

\begin{longtable}{@{}ll@{}}
\toprule
Objeto & Évaluation décimale (24 chiffres significatifs)\\
\midrule\endhead
$\alpha$ & $0.00729735256928380099728511\ldots$ \\
$H_5$ & $1.05457181764615446962582\ldots$ \\
$\eta_{\rm ret}$ & $1.05457181691503906113369\ldots$ \\
$A_{\deg}$ & $7.29735256928380099728511\ldots$ \\
$C^*_{\deg}$ & $2.12373833896864572426195\ldots$ \\
$q_+$ & $0.848377942576404374945737\ldots$ \\
$q_-$ & $0.913660151150557285295519\ldots$ \\
$r_+$ & $0.999999628548249363178695\ldots$ \\
$r_-$ & $0.999724137189518364131517\ldots$ \\
$\widehat\varepsilon$ & $0.999999257096636702760442\ldots$ \\
$\widehat\mu$ & $0.999448350479326935089982\ldots$ \\
$\widehat Z$ & $0.999724508538937215455554\ldots$ \\
$\widehat c$ & $1.00027631048615752610639\ldots$ \\
$\gamma$ & $0.274007672694459274747255\ldots$ \\
$G_{\rm Cat}$ & $0.915965594177219015054604\ldots$ \\
$\rho$ & $0.267949192431122706472554\ldots$ \\
\bottomrule
\end{longtable}


Le calcul s’effectue avec 120 chiffres de travail ; on montre 24 chiffres
significatifs pour faciliter la comparaison entre grandeurs de différents
ordres. Les 120 chiffres sont une précision d’évaluation, non une incertitude
expérimentale ni, à eux seuls, un certificat d’erreur par intervalles.
Le paquet conserve les entrées engendrées, leurs empreintes, les opérations et
les résidus des identités vérifiées. Les preuves algébriques du
corps du texte établissent ces identités indépendamment de ce tableau.

Le sens des quatre dernières coordonnées constitutives est conservé
lors de l’évaluation : $\widehat\varepsilon$ et $\widehat\mu$ sont les réponses
quadratiques par feuillet, $\widehat Z$ est leur orientation relative et
$\widehat c$ est l’inverse de leur produit. La proximité de certains chiffres avec
l’unité exprime le petit défaut du quotient $3$--$4$ sur ces
canaux. Son origine se voit exactement dans
\[
 1-f(s)=\frac{s^3}{1+s+s^2+s^3}.
\]
La formule permet d’étudier la réponse avant de lui attribuer une grandeur
dimensionnelle : le signe du défaut, son ordre cubique et sa monotonie sont
des conséquences du même opérateur.

\subsection{Coordinatisations d’unités et covariance}

Soient $\mathcal L_S$, $\mathcal L_Q$ et $\mathcal L_C$ les droites
d’action, de charge et de vitesse. Une coordinatisation physique positive évalue les sections
$u_S$, $u_Q$ et $u_C$ en unités d’action, de charge et de vitesse. La réalisation
\eqref{iii:eq:dimensions} produit alors une impédance, une vitesse,
une perméabilité et une permittivité, avec leurs dimensions respectives.
Les coordonnées du tableau précédent restent adimensionnelles.

\begin{proposition}[Covariance de la réalisation constitutive]
Si les sections de référence changent selon
$u_S'=a u_S$, $u_Q'=b u_Q$, $u_C'=d u_C$, avec $a,b,d>0$, les
coordonnées qui représentent les mêmes grandeurs satisfont
\[
 \widehat Z'=a^{-1}b^2\widehat Z,\qquad
 \widehat c'=d^{-1}\widehat c,\qquad
 \widehat\mu'=a^{-1}db^2\widehat\mu,\qquad
 \widehat\varepsilon'=adb^{-2}\widehat\varepsilon.
\]
Les relations $\widehat\mu'\widehat\varepsilon'=(\widehat c')^{-2}$
et $\widehat\mu'/\widehat\varepsilon'=(\widehat Z')^2$ sont conservées.
\end{proposition}
\begin{proof}
On égale chaque grandeur de \eqref{iii:eq:dimensions} à son expression
dans la nouvelle base et on isole sa coordonnée. Les facteurs du produit
sont $d^2$, et ceux du quotient sont $a^{-2}b^4$, exactement ceux requis
par les coordonnées transformées de $c$ et $Z$.
\end{proof}

Cette covariance distingue deux actes : construire une section et exprimer
cette section dans des unités choisies. Le théorème spectral fixe les coordonnées
dans la normalisation interne déclarée. Une comparaison SI doit également présenter
l’évaluation des trois sections de référence. Les ajuster
au moyen des grandeurs que l’on souhaite reproduire serait un étalonnage de
la coordinatisation, non une vérification indépendante de ces grandeurs.

Pour reconstruire les canaux à partir de coordonnées exprimées dans la coordinatisation
primée, on défait ce changement avant d’utiliser la racine cubique :
\begin{equation}
 \widehat Z=ab^{-2}\widehat Z',\quad
 \widehat c=d\widehat c',\quad
 \widehat\mu=ad^{-1}b^{-2}\widehat\mu',\quad
 \widehat\varepsilon=a^{-1}d^{-1}b^2\widehat\varepsilon'.
 \label{iii:eq:undo-units}
\end{equation}
Par exemple, les entrées de l’inverse spectral sont
$r_+=(\widehat Z\widehat c)^{-1/2}$ et
$r_-=(\widehat Z/\widehat c)^{1/2}$, calculées après
\eqref{iii:eq:undo-units}. Introduire directement les coordonnées primées
dans ces expressions changerait la normalisation de l’opérateur et pourrait
les faire sortir de l’intervalle $(3/4,1)$. La covariance dimensionnelle conserve
les grandeurs physiques ; l’inversion cubique reconstruit le transport
dans la coordinatisation interne dans laquelle il a été défini.
